\section{Introduction}\label{sec:intro}

\subsection{A list, and a question}

The foundations of mathematics are usually met as a list. Extensionality, empty set, pairing,
union, power set, separation, infinity, replacement, foundation, choice: ten statements about
membership \citep{Zermelo1908, Fraenkel1922} from which, we are told, the rest of mathematics
follows. The list is adopted because it works. It carries analysis, algebra and topology, and it
has not produced a contradiction.

This paper asks where the list comes from. Not whether it is true or consistent, but whether its
entries are ten independent stipulations, or whether some of them are consequences of something
simpler, with only a few genuinely \emph{chosen}. Our answer is that the list sorts into two
columns. Some axioms are \emph{free}: they hold automatically of anything produced by a modest
finite construction. Others are \emph{purchased}: they mark the exact points at which that
construction must be told to go on past where it would stop on its own. The free column is
proved; the purchased column is stated as a small set of named hypotheses, each with a theorem
or countermodel showing why it cannot be dropped.

\subsection{The idea in brief}

The organizing vocabulary comes from Six Birds Theory (\textsc{SBT})~\citep{SBTFoundationsI}, a
framework for describing how stable objects arise when a complicated substrate is observed
through a limited interface. We use only a small part of it, and we use it as an organizing
instrument, not as a source of new set theory. Three roles matter here.
\begin{itemize}[itemsep=2pt,topsep=3pt]
\item A \emph{lens} records what an observer can see. For sets, the substrate is the class of
finite membership graphs, and the lens records only the membership profile of the distinguished
node: which sets are its members, which sets are theirs, and so on. Node names and duplicated
structure are invisible.
\item A \emph{packaging map} replaces each graph by a canonical representative of what the lens
sees. For sets it collapses the part of the graph reachable from the distinguished node, as in
a finite Mostowski collapse, and re-presents it canonically; its objects are exactly the
hereditarily finite sets. Extensionality is the decision to identify graphs that the lens cannot tell apart.
\item An \emph{audit} is a certificate that strictly decreases along the structure. For sets it
is the rank, which decreases along membership. On finite graphs, Foundation says precisely that
such a certificate exists.
\end{itemize}
Two further patterns from the framework organize the rest of the paper. A fixed finite rule,
iterated, may keep producing new objects, but they stay of the same kind: everything it produces
is hereditarily finite, and it never leaves $\HF$. In this sense it \emph{saturates}. An object
of a new kind, such as an infinite set, comes only from \emph{changing the rule}, that is, from
adjoining something the old rule cannot produce. The free axioms are what the finite rules
already supply inside $\HF$; the purchased axioms are the named changes of rule. Figure~\ref{fig:overview} shows the resulting picture of the cumulative
hierarchy.

\begin{figure}[tbp]
\centering
\begin{tikzpicture}[
    font=\small,
    box/.style={draw, rounded corners=2pt, align=left, font=\footnotesize,
                inner sep=4pt, text width=5.4cm, execute at begin node=\raggedright},
    >={Stealth[length=2mm]}]
  % the cumulative hierarchy as a cone
  \coordinate (O) at (0,0);
  \coordinate (L) at (-2.4,6.0);
  \coordinate (R) at (2.4,6.0);
  \coordinate (Lw) at (-0.96,2.4);
  \coordinate (Rw) at (0.96,2.4);
  \fill[blue!10] (O) -- (Lw) -- (Rw) -- cycle;
  \fill[orange!12] (Lw) -- (L) -- (R) -- (Rw) -- cycle;
  \draw[thick] (L) -- (O) -- (R);
  \draw[densely dashed, thick] ($(Lw)+(-0.3,0)$) -- ($(Rw)+(0.3,0)$);
  \node[font=\footnotesize, anchor=west] at ($(Rw)+(0.32,0)$) {$\omega$};
  \foreach \t in {0.3,0.55,0.78} {
    \draw[gray!70] ($(O)!\t!(Lw)$) -- ($(O)!\t!(Rw)$);
  }
  \node[align=center, font=\footnotesize] at (0,1.25) {$\HF$};
  \node[align=center, font=\footnotesize] at (0,4.3) {stages $V_\alpha$\\$\alpha\ge\omega$};
  % left: free column
  \node[box, fill=blue!7, anchor=east] (F) at (-1.4,1.9) {%
    \textbf{Free} (Sections~\ref{sec:free}--\ref{sec:shadow})\\[2pt]
    Extensionality: the lens quotient\\
    Empty, Pairing, Union: packaging\\
    Separation: gates (subsets)\\
    Replacement: finite images\\
    Power set, Choice: finite forms\\
    Foundation: the rank audit};
  \draw[->, gray] (F.east) -- (-0.8,1.9);
  % right: purchased column
  \node[box, fill=orange!12, anchor=west] (B) at (2.3,4.4) {%
    \textbf{Purchased} (Section~\ref{sec:purchases})\\[2pt]
    \textbf{C-INF}: admit an inductive set\\
    \textbf{C-POW}: collect all subsets\\
    \textbf{C-REPL}: unbounded images\\
    \textbf{C-SEL}: sections of all surjections};
  \draw[->, gray] (B.west) -- (1.55,4.4);
  % barrier annotation
  \node[box, draw=none, anchor=west] (N) at (2.3,1.5) {%
    No finitely generated rule started\\in $\HF$ crosses the dashed line\\(Theorem~\ref{thm:nonattain})};
  \draw[->, gray] (N.west) -- (1.15,2.25);
\end{tikzpicture}
\caption{The free/purchased picture of the cumulative hierarchy. Below $\omega$ (blue), a finite
packaging construction supplies $\HF$, in which every axiom except Infinity holds; no finitely
generated rule iterated inside $\HF$ reaches the dashed line. Above it (orange), four named
commitments, on top of a structural premise that applies the free clauses at every stage
(Theorem~\ref{thm:shadow}), yield $\mathsf{ZF}$ and, with \textbf{C-SEL}, $\mathsf{ZFC}$.}
\label{fig:overview}
\end{figure}


\subsection{Main results}

We work first over the hereditarily finite sets $\HF$, where everything is finite and can be
checked, and then pass to the infinite by explicit hypotheses.

\paragraph{The free part (\S\ref{sec:free}).} On finite well-founded membership graphs, the
packaging map is idempotent and its fixed points correspond bijectively to $\HF$
(Theorem~\ref{thm:p5}). Closing the empty family under the finite packaging operations---empty
set, pairing, union, subsets and finite images---gives exactly $\HF$ (Theorem~\ref{thm:hfclosure}),
and $(\HF,\in)$ satisfies every axiom of $\mathsf{ZFC}$ except Infinity, each witnessed by a
named clause, by finite selection, or by the rank audit (Theorem~\ref{thm:hfsat}). On finite graphs, Foundation is
equivalent to the existence of a strictly decreasing rank function (Theorem~\ref{thm:tfae}).

\paragraph{Infinity is a change of rules (\S\ref{sec:infinity}).} No finitely generated
completion rule, iterated from hereditarily finite sets, ever produces an inductive set
(Theorem~\ref{thm:nonattain}). An infinite set has to be admitted; it cannot be reached.

\paragraph{Forgetting the stages is lossless (\S\ref{sec:shadow}).} The usual axioms mention
only membership, while the construction is staged. Over $\HF$ rank is first-order definable from
membership, and every sentence about the staged structure translates into an equivalent sentence
about membership alone (Theorem~\ref{thm:translation}).

\paragraph{The purchased part (\S\ref{sec:purchases}).} Infinity, Power set, Replacement and
Choice enter as four named commitments. Infinity is provably not free
(Theorem~\ref{thm:nonattain}). Choice is, over $\mathsf{ZF}$, exactly the statement that every
surjection of sets has a section (Theorem~\ref{thm:choice}); the particular membership lens needs
no choice, because its sections are given explicitly. Under a stated structural premise and the
four commitments, the membership structure satisfies $\mathsf{ZF}$, and with Choice $\mathsf{ZFC}$
(Theorem~\ref{thm:shadow}). Standard countermodels show that none of the four can be omitted
(Proposition~\ref{prop:nonredundancy}). Table~\ref{tab:factorization} records the resulting
ledger.

\paragraph{Forcing (\S\ref{sec:forcing}).} The Six Birds framework contains a counting
lemma---informally called the ``finite forcing lemma''~\citep{SBTFoundationsI}---stating that a
random predicate on $N$ points is constant on every block of a $K$-block partition with
probability $2^{-(N-K)}$. We show that this lemma counts the failures of the same disagreement
conditions that Cohen forcing uses~\citep{Cohen1963, Cohen1966}, and we compute exactly which
finite conditions have no extension meeting them (Theorem~\ref{thm:densitydefect}). Over
$\omega$ the failures disappear. As a consequence, a Cohen-generic real over a countable
transitive model $M$ lies outside $M$ and splits every infinite set coded in $M$
(Theorem~\ref{thm:novelty}), and adjoining it strictly enlarges the Boolean algebra of
$M$-coded predicates (Theorem~\ref{thm:strict}). The implication runs one way only: a real can
be new and split every infinite coded set without being generic
(Proposition~\ref{prop:converse}). The exact condition is that the real avoid every $M$-coded
closed nowhere dense set (Theorem~\ref{thm:recognition}). Cohen reals form a comeagre set of
measure zero, so the finite probability $1-2^{-(N-K)}$ describes novelty, not genericity
(\S\ref{ssec:category}). The continuum hypothesis appears as a question the ground model does
not settle (\S\ref{sec:ch}).

\paragraph{Guards and verification (\S\S\ref{sec:atlas}--\ref{sec:mech}).} Eight countermodels
block the most tempting overstatements, such as ``extensionality is forced'' or ``Foundation is
free without staging.'' Selected finite cores are verified in Lean~4 without external libraries.
An automatic audit, run on every build, checks that each of the 371 project declarations depends
on no axiom other than propositional extensionality and quotient soundness; in particular, the
finite lift used to price Choice is verified without Choice. Reproducible numerical exhibits
accompany the counting results.

\subsection{Scope and grades}

We do \emph{not} derive $\mathsf{ZF}$ or $\mathsf{ZFC}$ from finite principles: the infinitary
content enters only through named hypotheses. We prove no consistency statement beyond the
classical results we cite, and no new independence result; the forcing section reorganizes
classical facts and adds finite and structural statements about them. We do not claim that
$\mathsf{ZFC}$ is the unique or most natural set theory. The interpretive reading in
\S\ref{sec:discussion} comes after the theorems and is never used to prove them.

Because the paper mixes statements of different strength, every numbered result carries a tag.
\inhouse{} marks a result proved here without set-theoretic hypotheses beyond the ambient
metatheory; some of these have Lean-verified cores. \condg{} marks a result proved under named
hypotheses or cited classical imports. \modelg{} marks a statement about specific models obtained
from cited results. \guardg{} marks a countermodel. \metag{} marks an interpretive reading that is
deliberately not a theorem. Section~\ref{sec:scope} collects the boundaries in one place.
